\section{The theorem and its precise constructive meaning}\label{sec:statement}

\subsection{What is proved}

The usual Jordan theorem says that the complement of a simple closed plane
curve has exactly two connected components, one bounded and one unbounded,
and that the curve is their common boundary. In constructive arithmetic,
both the input and the conclusion must say what information is available.
In particular, a negative assertion that a point is not on the curve does
not supply a positive distance bound.

Our input is a quantitatively presented embedding
\[
 \gamma:\T\longrightarrow\R^2,\qquad \T=\R/\mathbb Z,
\]
where the circle distance is
\[
 d_\T(s,t)=\min(|s-t|,\,1-|s-t|)
 \quad(0\le s,t<1),
\]
and the plane uses the maximum norm. Rational coordinates, rational
parameters modulo one, and this norm make all finite computations exact.

\begin{definition}[Quantitative presentation]\label{def:presentation}
A presentation consists of a rational approximation routine
$G(q,n)\in\Q^2$, for $q\in\Q/\mathbb Z$, and nondecreasing functions
$U,V:\N\to\N$ with $U(k),V(k)\ge k+8$. Its associated curve has
\begin{align}
 \norm{G(q,n)-\gamma(q)}&\le 2^{-n},\label{eq:name}\\
 d_\T(s,t)\le 2^{-U(k)}
   &\ \Longrightarrow\
       \norm{\gamma(s)-\gamma(t)}\le 2^{-k},\label{eq:forward}\\
 d_\T(s,t)\ge 2^{-k}
   &\ \Longrightarrow\
       \norm{\gamma(s)-\gamma(t)}\ge 2^{-V(k)}.\label{eq:inverse}
\end{align}
All three routines are primitive recursive for an iPRA instance.
The second modulus is supplied data.
\end{definition}

The normalization of the moduli loses nothing: replace a modulus by the
maximum of its preceding values and $k+8$. The inverse condition is a
positive quantitative injectivity assertion. For rational $s,t$ it also
gives the usable implication
\[
 \norm{\gamma(s)-\gamma(t)}<2^{-V(k)}
 \ \Longrightarrow\ d_\T(s,t)<2^{-k};
\]
the comparison of the rational parameter distances is decidable.
The conditions extend to fast Cauchy parameter names by rational
approximation; the explicit padding used when the input assumptions
are stated only on rational parameters is given in
Section~\ref{sec:padding}.

Write $\Gamma$ for the image of the completed parameter circle.
An \emph{off-curve point} is a named point $x$ together with a natural
number $k_x$ certifying
\begin{equation}
 \norm{x-\gamma(t)}\ge 2^{-k_x}\quad\hbox{for every }t\in\T.
 \label{eq:off}
\end{equation}
The \emph{apartness complement} is the collection of points with such a
witness. A polygonal route has a finite list of vertices and a positive
rational lower bound for the distance of every one of its segments from
$\Gamma$. Its intermediate vertices can be rational even when its two
endpoints are named real points.

\begin{theorem}[Quantitative constructive Jordan theorem]\label{thm:main}
For a presentation as in Definition~\ref{def:presentation}, there are
primitive recursive procedures, relative to its fixed routines, with the
following properties.
\begin{enumerate}[label=\textnormal{(\arabic*)}]
\item \textbf{Classification.} Every off-curve point has a bit
$b_\Gamma(x)\in\{0,1\}$, independent of the chosen clearance certificate
and sufficiently accurate approximations. The bit is locally constant
on the apartness complement.
\item \textbf{Routing.} Given two off-curve points with equal bits, the
procedures return a polygonal route between them and a positive rational
clearance for that route.
\item \textbf{Separation.} A path with a supplied uniform continuity
modulus and positive clearance from $\Gamma$ has equal endpoint bits.
If instead its endpoints have different bits, then for every $n$ a
bounded calculation returns rational parameters $s,t$ such that
\[
       \norm{\alpha(s)-\gamma(t)}<2^{-n}.
\]
\item \textbf{Both sides.} Given a named parameter $t$ and $n$, the
procedures return rational points $a_0,a_1$ with explicit positive
clearances, with
\[
 b_\Gamma(a_0)=0,\qquad b_\Gamma(a_1)=1,\qquad
 \norm{a_i-\gamma(t)}<2^{-n}.
\]
\item \textbf{Size of the regions.} A primitive recursive bound $R$
for $\Gamma$ can be calculated. Every point with
$\norm{x}>R$ has bit zero. Thus the bit-one region is bounded and
inhabited, and the bit-zero region is unbounded and inhabited.
\item \textbf{Common boundary.} The image $\Gamma$ is exactly the set
of named points every neighborhood of which contains a rational
off-curve point of each bit. In the reverse direction a primitive recursive
circle name for a preimage is reconstructed from the point name.
\end{enumerate}
Consequently the apartness complement has exactly two classes under
the relation of being joined by a polygonal route with positive
clearance. Each class is open in the quantitative sense, and every
curve point is approached from both classes.
\end{theorem}

The term \emph{component} in this constructive conclusion has the explicitly
stated path-and-clearance meaning. We do not infer a uniform lower bound
from a merely pointwise claim that a continuous path misses the curve.
Nor do we replace the approximate intersection conclusion by the choice
of an exact zero or an exact crossing parameter. Those are different
selection assertions.

There is no rectifiability, differentiability, curvature, polygonality,
or supplied extension to a plane homeomorphism assumption. Classically,
continuous circle embeddings have forward and inverse moduli by
compactness. Here the moduli are part of the constructive presentation;
for the first-order iPRA claim they must be primitive recursive.

\subsection{How a first-order arithmetic theory speaks about these objects}
\label{sec:names}

iPRA has natural-number variables, function symbols for primitive recursive
functions, intuitionistic first-order logic, and quantifier-free induction.
It does not have a sort containing all real functions.
The preceding notation is an economical way to describe an
\emph{arithmetical theorem schema}.

Fix a finite collection of primitive recursive definitions for $G,U_0,V_0$
and, where needed, point and path names, possibly with number parameters
$\vec a$. Signed integers, positive-denominator rational pairs, finite
lists, and pairs of coordinates are coded by natural numbers.
All rational comparisons become decidable integer comparisons.
The public moduli $U,V$ are obtained by the explicit padding in
Section~\ref{sec:padding}. The assertion that the raw presentation
is valid can be written using the universal closures of the
following rational inequalities:
\begin{align}
\norm{G(s,n)-G(s,m)}
 &\le 2^{-n}+2^{-m},\label{eq:arithC}\\
d_\T(s,t)\le 2^{-U_0(k)}
 &\ \Longrightarrow\
 \norm{G(s,n)-G(t,n)}\le 2^{-k}+2^{1-n},\label{eq:arithU}\\
d_\T(s,t)\ge 2^{-k}
 &\ \Longrightarrow\
 \norm{G(s,n)-G(t,n)}\ge 2^{-V_0(k)}-2^{1-n}.
 \label{eq:arithV}
\end{align}
Variables $s,t$ range over codes of rational residues; invalid codes can
be sent to a fixed default. A rational name for a point satisfies the
analogous Cauchy inequalities. For a point name $X(n)$, condition
\eqref{eq:off} is expressed by
\[
 \forall q,n\quad
 \norm{X(n)-G(q,n)}\ge 2^{-k_x}-2^{1-n}.
\]
Only natural-number quantifiers occur. The limiting notation in
\eqref{eq:name}--\eqref{eq:inverse} abbreviates these inequalities and
their consequences at arbitrary requested precision.

For each fixed tuple of definitions, our conclusion is an iPRA proof of
the corresponding first-order implications from these regularity
premises to the correctness of explicitly described primitive recursive
outputs. We do not assume a primitive recursive universal evaluator
for arbitrary primitive recursive program codes. There is no such
uniform shortcut in this proof.

The distinction between a name and a real-function quantifier is
substantive. For example, a primitive recursive rational approximation
routine can have a useful continuity modulus whose computational
complexity is not primitive recursive. The theorem does not silently
manufacture the missing primitive recursive modulus.

\subsection{Finite nets and rational distance calculations}
\label{sec:nets}

For an integer $\ell$, set
\begin{equation}
 N(\ell)=2^{\,U(\ell)+1},\qquad
 t_j=j/N(\ell),\qquad Y_j=G(t_j,\ell).
 \label{eq:net}
\end{equation}
Every curve point is within $2^{1-\ell}$ of some $Y_j$:
the parameter mesh costs at most $2^{-\ell}$ and evaluation costs
another $2^{-\ell}$. Every $Y_j$ is within $2^{-\ell}$ of the
corresponding curve point. Both directions are proved by choosing the
appropriate mesh interval and applying \eqref{eq:arithU}.
The same construction works on a rational parameter interval, including
either half of the circle.

For a rational segment $L=[a,b]$, its distance to a rational point $y$
in the maximum norm is rational and is found by a bounded calculation.
Indeed, put
\[
 f(t)=\max\{|a_x+t(b_x-a_x)-y_x|,
             |a_y+t(b_y-a_y)-y_y|\},\qquad 0\le t\le1.
\]
It is the maximum of four affine functions. Its minimum is attained at
an endpoint or an intersection of two of the four affine functions;
retain only intersections with nonzero denominator and parameter in
$[0,1]$. There are at most eight candidates, including the two
endpoints. On each interval between consecutive candidates one affine
function gives the maximum, so its minimum is at an endpoint (or it is
constant). This proves the formula without an extremum theorem.

Consequently
\[
 D_L=\min_j\dist(L,Y_j)
\]
is a rational number. If the net error is at most $\rho$, then
\begin{equation}
  \big|D_L-\dist(L,\Gamma)\big|\le\rho.
  \label{eq:distance-net}
\end{equation}
One need not introduce an infimum as a new object: the two inequalities
mean that every segment--curve pair is at least $D_L-\rho$ apart, and
that a boundedly selected mesh parameter and segment parameter give
distance at most $D_L+\rho$. The same observation applies to a point
or a finite polygonal route.

A curve bound is also available by a finite calculation. For example,
sample with $N(0)$ and precision $2$, and take
\begin{equation}
 R=3+\max_{j<N(0)}\ceil{\norm{G(j/N(0),2)}}.
 \label{eq:R}
\end{equation}
Then $\Gamma\subset[-R,R]^2$. These nets and distance estimates are the
only approximation apparatus used below.

\subsection{Origin of the method}

The continuous constructive routing and local-boundary strategy is due
to Berg, Julian, Mines, and Richman \cite{BJMR}. Our presentation develops
its finite arithmetic content: a rational affine honeycomb, parity in
place of angles, a cycle-edge multiplicity proof of factorization,
explicit dyadic scales, and a symmetric chain calculation for the
boundary argument. The broader constructive setting is discussed by
Bridges \cite{Bridges}. Nguyen and Cook \cite{NC} give an independent
weak-arithmetic treatment of discrete Jordan principles; their theorem
is context, not an unproved lemma imported into the argument here.
